\section{Introduction}

\subsection{About nlist}

\noindent nlist is a self-hosted file listing and file management platform written in Go. Its main purpose is to expose many different storage backends through one central interface. Users can access files through the web UI and the HTTP API, while administrators can manage users, permissions, and storage configuration.

\noindent The backend is built with Gin and exposes much of its application logic through routes under \texttt{/api}. In addition to the browser-based interface, nlist also supports protocols and services such as WebDAV, S3-compatible access, FTP, and SFTP. This makes it a general-purpose file access gateway for different clients and environments.

\noindent Internally, nlist separates routing, authentication, and storage-related operations across different packages. Users can authenticate through endpoints such as \texttt{/api/auth/login}, receive a token, and then access protected routes for file operations and administration. Because these components control access to potentially large amounts of data, the authentication layer is a particularly important security boundary.

\subsection{Security Issue}

\noindent nlist tries to counter brute-force attacks to the auth API endpoint by counting failed login attempts per client IP. After 5 failed attempts, the IP is blocked for 5 minutes and further requests terminate in an HTTP 429 error. At first this seems like a reasonable rate-limiting mechanism.

\noindent However, the implementation uses \texttt{c.ClientIP()} as the key for this protection. The server is created with \texttt{gin.New()} and no explicit trusted proxy configuration. So Gin can derive the client IP from forwarding headers such as \texttt{X-Forwarded-For}. This implementation would be fine if the application were deployed behind a configured reverse proxy and only known addresses were trusted. But in this case an attacker can send every login request with a different fake \texttt{X-Forwarded-For} value:

\noindent All of these requests are treated as different clients by the application although they come from the same machine. The implemented rate-limiting security mechanism is never triggered and as a result the attacker can send unlimited login attempts to brute-force the login credentials.